\documentclass[11pt]{article}
\usepackage[margin=1in]{geometry}
\usepackage{amsmath,amssymb,amsthm,mathtools}
\usepackage{booktabs}
\usepackage{microtype}
\usepackage[hidelinks]{hyperref}
\usepackage{enumitem}

\newtheorem{theorem}{Theorem}
\newtheorem{lemma}[theorem]{Lemma}
\newtheorem{corollary}[theorem]{Corollary}
\newtheorem{proposition}[theorem]{Proposition}
\theoremstyle{definition}
\newtheorem{definition}[theorem]{Definition}
\newtheorem{remark}[theorem]{Remark}

\newcommand{\Z}{\mathbb Z}
\newcommand{\F}{\mathbb F}
\newcommand{\M}{M_8}
\newcommand{\tr}{\operatorname{tr}}

\title{\textbf{An Eight-State Operator and Transpose Geometry over Dyadic Rings}\\
Trace-Zero Modules, Mixed-Radix Addressing, and Tensor-Lift Kernels}
\author{Muhammad Arshad}
\date{September 2026}

\begin{document}
\maketitle

\begin{abstract}
We give an exact algebraic synthesis of the 8-state operator geometry used by the Quantinion/MPRC construction over the dyadic ring $R_m=\Z/2^m\Z$, with the specialization $R_8=\Z_{256}$. The eight states are separated into one reference state and seven non-reference states. The quadratic-residue set $\{1,2,4\}\subset\Z_7$ canonically orients the $21$ unordered pairs among the seven non-reference states and recovers the classical Frobenius group $F_{21}$ generated by $x\mapsto x+1$ and $x\mapsto2x$. Inside the full $8\times8$ operator module, the trace-zero submodule is free of rank $63$ and admits an explicit decomposition $63=14+21+28=56+7$. Fourteen reference couplings generate the whole trace-zero module under commutators. We then connect this operator space to the frozen QH4 active-position formula: after writing $\gamma-1=9\theta+3a+p$, the address becomes $m=64\theta+21a+7p+\sigma$, followed by the unit stride $z=7m\bmod256$. This exhibits the active ring as four sectors of $63$ non-vacuum addresses. A local Arshad Transpose $T_k(a,p)=(p+k,a-k)$ is an involution preserving $a+p$ and transporting $p-a$ by $j\mapsto-(j+2k)$; on the $252$ active states it has $84$ fixed states and $84$ two-cycles. Finally, for the canonical $\ell$-fold tensor action
\[
\rho_\ell(A)=\sum_{r=1}^{\ell}I^{\otimes(r-1)}\otimes A\otimes I^{\otimes(\ell-r)}
\]
we prove
\[
\ker\rho_\ell=\{\lambda I_8:\ell\lambda=0\pmod{2^m}\}.
\]
For $\Z_{256}$ this gives kernels $\{0,128I\}$ at order two, $\{0\}$ at order three, and $\{0,64I,128I,192I\}$ at order four. The group $F_{21}$ and elementary matrix-unit commutator identities are classical; the contribution is the exact coupling of these structures to the QH4 address, the Transpose orbit decomposition, and the dyadic tensor-lift kernel.
\end{abstract}

\noindent\textbf{Keywords:} matrices over rings; trace-zero matrices; Frobenius group; tensor products; dyadic rings; matrix commutators; finite operator geometry.

\medskip
\noindent\textbf{Code and reproducibility:} A pure-integer verifier accompanies the public preprint repository and checks the finite group, operator basis, address bijection, transpose orbits, and tensor-lift kernels.

\section{Introduction}
The numerical sequence
\[
7,\quad21,\quad63,\quad64,\quad252,\quad256
\]
appears repeatedly in the Quantinion/MPRC construction. A central objective of this paper is to separate identities that are merely arithmetic from identities that have exact operator meaning.

The starting state geometry is the binary generator
\[
g=(1,2,4).
\]
Its three binary coordinates generate eight subset states. The full endomorphism module on eight states has $8^2=64$ matrix units. Once one state is chosen as a reference, the seven remaining states admit a natural orientation based on the residue set $\{1,2,4\}$ modulo seven. The same orientation is classical: the affine transformations $x\mapsto ax+b$, with $a\in\{1,2,4\}$ and $b\in\Z_7$, form the Frobenius group $F_{21}$ of order $21$ \cite{CostaPavone2024}. We do not claim novelty for this group.

The first new task is instead to understand how the seven-state orientation sits inside the $8\times8$ operator module over a non-field ring. We show that the trace-zero submodule has an explicit free basis of rank $63$, that the orientation yields the exact module split
\[
63=14+21+28,
\]
and that only the $14$ reference couplings are needed to generate the whole trace-zero module under commutators.

The second task is to connect the operator counts to the pre-existing QH4 active-position address rather than to impose a new coordinate system. The frozen address contains $36=4\cdot9$ layers and seven $\sigma$ states. Rewriting its $\gamma$ coordinate reveals the coefficients $64$, $21$, $7$, and $1$ exactly. A local Transpose on the $3\times3$ subcoordinate then supplies a rigorous $21+21+21$ state-space partition inside each $63$-state sector.

The third task concerns composition. We prove a tensor-lift kernel theorem over $\Z/2^m\Z$. It explains precisely which central scalar states disappear at even tensor orders and why the third-order $8^3=512$ action is faithful over $\Z_{256}$.

\paragraph{Contributions.}
\begin{enumerate}[leftmargin=*]
\item We state explicitly the classical $F_{21}$ orientation on seven points and separate it from the new dyadic-ring results.
\item We construct a free rank-$63$ trace-zero submodule of $M_8(\Z/2^m\Z)$ and decompose it as $14+21+28$ according to reference couplings, oriented non-reference edges, reverse edges, and diagonal differentials.
\item We prove that the $14$ reference couplings generate all $63$ trace-zero basis directions by elementary commutators.
\item We derive the mixed-radix identity
\[
m=64\theta+21a+7p+\sigma,\qquad z=7m\pmod{256},
\]
directly from the frozen QH4 address and prove exact $252$-state coverage.
\item We prove the local Transpose involution and its $84$ fixed plus $84$ two-cycle global orbit decomposition.
\item We prove the tensor-lift kernel formula $\ker\rho_\ell=\{\lambda I_8:\ell\lambda=0\}$ over every dyadic ring.
\end{enumerate}

\paragraph{Novelty boundary.}
The Frobenius group of order $21$ and oriented Fano-plane realizations are classical and have recent explicit treatments \cite{CostaPavone2024}. Trace-zero matrices and matrix commutators over commutative rings are also a classical algebraic topic; see, for example, Suwama \cite{Suwama2023}. This paper does not claim either of those results as new. The proposed contribution is the exact synthesis with the Quantinion/QH4 address, the specific Transpose orbit law, the explicit $14$-generator realization of the trace-zero module in this state labeling, and the dyadic tensor-lift kernel theorem.

\section{The eight-state operator space}
Let
\[
R_m=\Z/2^m\Z.
\]
Index the eight states by one reference symbol $\infty$ together with the seven elements of $\Z_7$:
\[
\Omega=\{\infty\}\sqcup\Z_7.
\]
Let $E_{ab}$ denote the $8\times8$ matrix unit with a one in row $a$, column $b$, and zeros elsewhere.

The full operator module $M_8(R_m)$ is free of rank $64$, with basis
\[
\{E_{ab}:a,b\in\Omega\}.
\]

The trace map
\[
\tr:M_8(R_m)\to R_m
\]
is surjective because $\tr(E_{\infty\infty})=1$. Hence
\[
M_8(R_m)=\ker(\tr)\oplus R_m E_{\infty\infty}.
\]

\begin{proposition}[Trace-zero module]\label{prop:tracezero}
The trace-zero module
\[
\mathfrak T_8(R_m)=\ker(\tr)
\]
is free of rank $63$. An explicit basis is
\[
\{E_{ab}:a\ne b\}
\cup
\{H_a=E_{aa}-E_{\infty\infty}:a\in\Z_7\}.
\]
\end{proposition}

\begin{proof}
There are $56$ off-diagonal matrix units and seven diagonal differences, for a total of $63$. Every trace-zero matrix is obtained uniquely by taking its off-diagonal entries directly and expressing its diagonal relative to the $\infty$ entry. Linear independence follows by inspection of the off-diagonal coordinates and then of the seven non-reference diagonal coordinates.
\end{proof}

Thus
\[
\operatorname{rank}_{R_m}\mathfrak T_8(R_m)=56+7=63.
\]

\section{The classical $21$-orientation on seven points}
Let
\[
Q=\{1,2,4\}\subset\Z_7^\times.
\]
Its negative set is
\[
-Q=\{6,5,3\},
\]
and
\[
Q\cap(-Q)=\varnothing,\qquad Q\cup(-Q)=\Z_7\setminus\{0\}.
\]

For an unordered pair $\{x,y\}$ of distinct points in $\Z_7$, orient it as
\[
x\longrightarrow y
\quad\Longleftrightarrow\quad
y-x\in Q.
\]

\begin{proposition}[Unique pair orientation]\label{prop:orientation}
Every unordered pair in $\Z_7$ receives exactly one orientation. Consequently there are exactly
\[
7\cdot3=\binom72=21
\]
oriented representatives.
\end{proposition}

\begin{proof}
For $x\ne y$, the two directed differences are $d=y-x$ and $-d=x-y$. Exactly one belongs to $Q$ because $Q$ and $-Q$ partition the six nonzero residues.
\end{proof}

Multiplication by two permutes $Q$ cyclically:
\[
1\mapsto2\mapsto4\mapsto1.
\]
The affine maps
\[
x\mapsto ax+b,\qquad a\in Q,\quad b\in\Z_7,
\]
form the classical Frobenius group $F_{21}$ \cite{CostaPavone2024}. The integer pair products of the generator $(1,2,4)$ are
\[
(1\cdot2,1\cdot4,2\cdot4)=(2,4,8),
\]
which reduce modulo seven to $(2,4,1)$, the same order-three cycle.

\section{A $63=14+21+28$ operator decomposition}
Separate the trace-zero basis into three families.

First, the reference couplings:
\[
\mathcal R=
\{E_{a\infty},E_{\infty a}:a\in\Z_7\},
\qquad |\mathcal R|=14.
\]

Second, the $21$ canonically oriented non-reference edges:
\[
\mathcal F=
\{E_{xy}:x,y\in\Z_7,\ y-x\in Q\},
\qquad |\mathcal F|=21.
\]

Third, the reverse edges together with the diagonal differences:
\[
\mathcal B=
\{E_{yx}:E_{xy}\in\mathcal F\}
\cup
\{H_a:a\in\Z_7\},
\]
so
\[
|\mathcal B|=21+7=28.
\]

\begin{theorem}[Operator-space decomposition]\label{thm:decomp}
There is a direct sum of free $R_m$-modules
\[
\mathfrak T_8(R_m)
=
\langle\mathcal R\rangle
\oplus
\langle\mathcal F\rangle
\oplus
\langle\mathcal B\rangle,
\]
with ranks
\[
63=14+21+28.
\]
\end{theorem}

\begin{proof}
The three families are disjoint subsets of the explicit basis in Proposition~\ref{prop:tracezero}, and together they exhaust that basis.
\end{proof}

\begin{remark}
This is a module decomposition. We do not claim that the three summands define a Lie grading or a decomposition into ideals.
\end{remark}

\section{Fourteen reference couplings generate the trace-zero module}
For matrix units,
\[
E_{ab}E_{cd}=\delta_{bc}E_{ad}.
\]

\begin{theorem}[Fourteen-generator commutator closure]\label{thm:14gen}
The $R_m$-module generated by $\mathcal R$ and its commutators contains the entire trace-zero module $\mathfrak T_8(R_m)$.
\end{theorem}

\begin{proof}
For distinct $a,b\in\Z_7$,
\[
[E_{a\infty},E_{\infty b}]
=
E_{ab}.
\]
For $a=b$,
\[
[E_{a\infty},E_{\infty a}]
=
E_{aa}-E_{\infty\infty}
=
H_a.
\]
Thus all $42$ non-reference off-diagonal units and all seven diagonal differences are generated. Together with the original fourteen reference couplings, these are exactly the basis of Proposition~\ref{prop:tracezero}.
\end{proof}

No division or field assumption is used; the result holds over every $R_m$.

\section{The frozen QH4 address in mixed-radix form}
The active QH4 address is taken as
\[
z
=
7\left(
7(\gamma-1)+\sigma+
\left\lfloor\frac{\gamma-1}{9}\right\rfloor
\right)
\pmod{256},
\]
where
\[
1\le\gamma\le36,\qquad1\le\sigma\le7.
\]

Write
\[
\gamma-1=9\theta+3a+p,
\]
with
\[
\theta\in\{0,1,2,3\},\qquad a,p\in\{0,1,2\}.
\]
Then
\[
\left\lfloor\frac{\gamma-1}{9}\right\rfloor=\theta,
\]
and therefore
\begin{align*}
7(\gamma-1)+\sigma+\left\lfloor\frac{\gamma-1}{9}\right\rfloor
&=
7(9\theta+3a+p)+\sigma+\theta\\
&=
64\theta+21a+7p+\sigma.
\end{align*}

Define
\[
m(\theta,a,p,\sigma)=64\theta+21a+7p+\sigma.
\]

\begin{theorem}[QH4 mixed-radix active address]\label{thm:qaddr}
For each fixed $\theta$, the values
\[
21a+7p+\sigma,
\qquad
a,p\in\{0,1,2\},\quad1\le\sigma\le7,
\]
are exactly the integers $1,\ldots,63$. Consequently $m$ runs through all elements of $\{0,\ldots,255\}$ except
\[
\{0,64,128,192\}.
\]
Since $\gcd(7,256)=1$, multiplication by seven permutes the ring, so $z=7m\bmod256$ also visits exactly $252$ positions.
\end{theorem}

\begin{proof}
For fixed $a$, $7p+\sigma$ runs through the interval $1,\ldots,21$. Adding $21a$ gives the consecutive intervals $1\ldots21$, $22\ldots42$, and $43\ldots63$. Adding $64\theta$ places these intervals in the four sectors. The only omitted value in each sector is its zero offset. The final stride is bijective because seven is a unit modulo $256$.
\end{proof}

Thus
\[
252=4\cdot63=4\cdot7\cdot9,
\qquad
256=4(63+1).
\]

\section{Local Arshad Transpose}
On the local $3\times3$ coordinate $(a,p)\in\Z_3^2$, define
\[
T_k(a,p)=(p+k,\ a-k)\pmod3,
\qquad k\in\Z_3.
\]

\begin{theorem}[Transpose involution and transported coordinate]\label{thm:transpose}
For each $k\in\Z_3$:
\begin{enumerate}[leftmargin=*]
\item $T_k^2=I$;
\item the sum coordinate $s=a+p$ is invariant;
\item the difference coordinate $j=p-a$ transforms as
\[
j\mapsto-(j+2k)\pmod3;
\]
\item the fixed points are exactly those with $j=-k$.
\end{enumerate}
\end{theorem}

\begin{proof}
Apply the definition twice:
\[
T_k(p+k,a-k)=(a-k+k,p+k-k)=(a,p).
\]
For the sum,
\[
(p+k)+(a-k)=a+p.
\]
For the difference,
\[
(a-k)-(p+k)=-(p-a+2k)=-(j+2k).
\]
Finally a point is fixed exactly when $p+k=a$, equivalently $p-a=-k$.
\end{proof}

There are three fixed local points, one for each value of the invariant sum $s$, and the remaining six points form three two-cycles.

\begin{corollary}[Global active-state orbit count]\label{cor:orbits}
For each $k$, extending $T_k$ independently across the seven $\sigma$ values and four $\theta$ sectors gives
\[
84
\]
fixed active states and
\[
84
\]
two-cycles. Thus
\[
84+2\cdot84=252.
\]
Within each $63$-state quadrant, the invariant $s\in\Z_3$ partitions the states into three disjoint $21$-state sectors.
\end{corollary}

\section{Canonical tensor lift}
Let $V=R_m^8$. For $A\in\operatorname{End}_{R_m}(V)$ and $\ell\ge1$, define
\[
\rho_\ell(A)
=
\sum_{r=1}^{\ell}
I^{\otimes(r-1)}
\otimes A
\otimes I^{\otimes(\ell-r)}
\in
\operatorname{End}_{R_m}(V^{\otimes\ell}).
\]

\begin{theorem}[Tensor-lift kernel]\label{thm:kernel}
For every $\ell\ge1$,
\[
\boxed{
\ker\rho_\ell
=
\{\lambda I_8:\ell\lambda=0\text{ in }R_m\}.}
\]
In particular,
\[
|\ker\rho_\ell|=\gcd(\ell,2^m).
\]
\end{theorem}

\begin{proof}
Write $A=(a_{ij})$. If some off-diagonal coefficient $a_{ij}$ with $i\ne j$ is nonzero, apply $\rho_\ell(A)$ to a pure tensor basis vector with all slots fixed and inspect the basis output in which exactly one chosen slot changes from $j$ to $i$. That output can arise only from the term in which $A$ acts on that slot. Hence no nonzero off-diagonal coefficient can belong to the kernel. Thus every kernel element is diagonal:
\[
A=\operatorname{diag}(a_0,\ldots,a_7).
\]

On the pure tensor basis vector
\[
e_{i_1}\otimes\cdots\otimes e_{i_\ell},
\]
the action is multiplication by
\[
a_{i_1}+\cdots+a_{i_\ell}.
\]
Compare two basis vectors that differ in exactly one slot, with values $i$ and $j$ there. Kernel membership forces
\[
a_i-a_j=0,
\]
so all diagonal entries are equal:
\[
A=\lambda I_8.
\]
The action then becomes
\[
\rho_\ell(\lambda I_8)=\ell\lambda I_{8^\ell}.
\]
Thus $\lambda I_8$ lies in the kernel exactly when $\ell\lambda=0$ in $R_m$. The number of solutions of $\ell\lambda=0$ modulo $2^m$ is $\gcd(\ell,2^m)$.
\end{proof}

\section{The $\Z_{256}$ composition ladder}
Set $m=8$. Then
\[
|\ker\rho_\ell|=\gcd(\ell,256).
\]
In particular,
\[
\ker\rho_2=\{0,128I\},
\]
\[
\ker\rho_3=\{0\},
\]
and
\[
\ker\rho_4=\{0,64I,128I,192I\}.
\]

The pair-level state dimension is $8^2=64$, while the cubic state dimension is $8^3=512$. The theorem shows a genuine difference in faithfulness: second order loses the central half-turn, whereas third order is faithful. At fourth order the central kernel has four scalar positions spaced by $64$, exactly the same spacing as the four distinguished QH4 vacuum addresses.

\begin{remark}
The last equality is an exact algebraic coincidence of subsets in $\Z_{256}$ under the chosen scalar embedding. No claim is made that it proves a physical identification between operator kernels and vacuum states.
\end{remark}

\section{Computational verification}
The reference program \texttt{verify\_quantinion\_operator\_geometry.py} uses exact integer and bit arithmetic only. It verifies:
\begin{enumerate}[leftmargin=*]
\item all $21$ affine maps of $F_{21}$ and closure under composition;
\item the unique orientation of the $21$ unordered pairs;
\item the $63$-element trace-zero basis and its binary rank;
\item all commutator identities in Theorem~\ref{thm:14gen};
\item the full $252/252$ mixed-radix QH4 address with exactly four omitted vacua;
\item the Transpose involution, invariant, transported coordinate, and $84+84$ global orbit counts;
\item scalar tensor-lift kernels for $\ell=1,\ldots,8$.
\end{enumerate}

\section{Relation to prior work}
The order-$21$ affine group used in Section 3 is the classical Frobenius group $F_{21}$. Costa and Pavone give oriented Fano-plane and $K_7$ realizations and explicitly study this group geometrically \cite{CostaPavone2024}. We use the group only as the canonical orientation algebra of the seven non-reference states.

For matrices over a commutative ring, every commutator has trace zero; the converse problem is subtler over general rings. Suwama studies precisely this trace-zero/commutator question and gives examples where trace-zero matrices fail to be single commutators \cite{Suwama2023}. The present paper does not claim that every trace-zero element of $M_8(R_m)$ is one commutator. Theorem~\ref{thm:14gen} states the weaker and explicit fact needed here: the trace-zero module is generated by the fourteen reference couplings together with iterated commutators and $R_m$-linear combinations.

The mixed-radix QH4 identity and the Transpose orbit law are specific to the frozen Quantinion/MPRC address used here. The tensor-lift kernel theorem is stated independently over dyadic rings.

\section{Limitations and open questions}
The operator decomposition is not identified with $\mathfrak g_2$, $\mathfrak{so}(8)$, the octonions, or any other classical Lie algebra. Earlier numerical similarities such as $7$, $14$, and $28$ are not used as proofs.

The Transpose theorem is proved for the stated local $3\times3$ involution and its product with the frozen $\theta$ and $\sigma$ coordinates. Extending the same invariant structure to other dimensions requires a separate definition and proof.

The tensor-lift theorem concerns the canonical derivation-like action $\rho_\ell$. Other tensor actions can have different kernels.

Finally, the $64/512$ faithfulness result is algebraic. Its value for AI kernels or physical modeling requires separate experiments and is not part of the theorem.

\section{Conclusion}
The Quantinion 8-state construction admits an exact operator interpretation. Seven non-reference states carry a classical $21$-edge orientation. Inside the $8\times8$ dyadic operator space, the trace-zero module has free rank $63$ and is generated from only fourteen reference couplings by elementary commutators. The frozen QH4 address exposes the same arithmetic through the exact mixed-radix law
\[
64\theta+21a+7p+\sigma,
\]
whose four sectors contain $63$ active positions each. The local Transpose is an involution that turns the $3\times3$ coordinate into three invariant $21$-state sectors per quadrant. Finally, the tensor-lift kernel theorem explains exactly which central scalars are lost at each composition order.

The resulting picture is not a new discovery of the classical $F_{21}$ group. Its contribution is the rigorous assembly of classical and new ring-native pieces into one explicit dyadic operator geometry.

\begin{thebibliography}{9}
\bibitem{CostaPavone2024}
S. Costa and M. Pavone,
"Orthogonal and oriented Fano planes, triangular embeddings of $K_7$, and geometrical representations of the Frobenius group $F_{21}$,"
\emph{AIMS Mathematics}, vol. 9, no. 12, pp. 35274--35292, 2024.
DOI: 10.3934/math.20241676.

\bibitem{Suwama2023}
M. Suwama,
"On trace zero matrices and commutators,"
\emph{Journal of Algebra}, vol. 616, pp. 26--48, 2023.
DOI: 10.1016/j.jalgebra.2022.11.005.
\end{thebibliography}

\end{document}
